% ===== Section 1: "My IQ is 270 ... or some other big number" (~16 min) =====
\topicopener{numbers}{IQ}{Իմ IQ-ն 270 է... կամ մի ուրիշ մեծ թիվ}{Ինչու է խելքը չափելն ավելի դժվար, քան թվում է}{confused}

% Asking "am I intelligent?" already proves it
\begin{frame}{Խելացի՞ ես}
  \begin{columns}[c]
    \begin{column}{0.58\textwidth}
      Եթե քեզ այդ հարցն ես տալիս, պատասխանը «այո» է։ Որովհետև դու...
      \begin{itemize}
        \item<2-> հասկանում ես «խելք» գաղափարը, որը ոչ հստակ սահմանում ունի, ոչ ֆիզիկական տեսք
        \item<3-> գիտակցում ես քեզ որպես առանձին անհատ
        \item<4-> համեմատում ես քեզ ուրիշների ու ինչ-որ իդեալի հետ
      \end{itemize}
      \onslide<5->{Ոչ մի այլ կենդանի սա չի կարողանում։\par}
    \end{column}
    \begin{column}{0.36\textwidth}
      \centering
      \mascot[1.9cm]{thinking}\par\smallskip
      \onslide<6->{%
      \begin{bookquote}
        Վատ չէ մի բանի համար, որն ըստ էության թեթև նևրոզ է։
      \end{bookquote}}
    \end{column}
  \end{columns}
\end{frame}

% Three classic definitions
\begin{frame}{Երեք սահմանում}
  \relax{\small Խելքը, ինչպես \textbf{հեգնանքը} կամ \textbf{ամառային ժամը}, բոլորը մոտավորապես հասկանում են, բայց քչերը կբացատրեն։}\par\smallskip
  \begin{columns}[T]
    \begin{column}{0.31\textwidth}
      \begin{card}[colback=peach!35!cardwhite]
        {\blackfont Բինե և Սիմոն}\par
        {\scriptsize առաջին IQ թեստերից մեկը}\par\smallskip
        «Լավ դատել, լավ հասկանալ, լավ տրամաբանել. սրանք են խելքի գլխավոր գործողությունները»
      \end{card}
    \end{column}
    \begin{column}{0.31\textwidth}
      \onslide<2->{%
      \begin{card}[colback=mint!35!cardwhite]
        {\blackfont Դեյվիդ Վեքսլեր}\par
        {\scriptsize WAIS թեստը՝ մինչ օրս}\par\smallskip
        «Նպատակային գործելու ու միջավայրի հետ արդյունավետ գլուխ հանելու կարողություն»
      \end{card}}
    \end{column}
    \begin{column}{0.31\textwidth}
      \onslide<3->{%
      \begin{card}[colback=lavender!40!cardwhite]
        {\blackfont Ֆիլիպ Վերնոն}\par
        {\scriptsize ոլորտի մեկ այլ մեծ անուն}\par\smallskip
        «Հասկանալու, կապեր տեսնելու ու տրամաբանելու ընդհանուր կարողություն»
      \end{card}}
    \end{column}
  \end{columns}
\end{frame}

% What everyone agrees on, and the padlocked gate
\begin{frame}{Ինչում են բոլորը համաձայն}
  \begin{columns}[c]
    \begin{column}{0.45\textwidth}
      Խելքը ուղեղի կարողությունն է անելու... \textbf{բաներ}։\par\medskip
      \onslide<2->{Ավելի ճիշտ՝ \textbf{տեղեկություն մշակելու ու օգտագործելու} կարողությունը. տրամաբանություն, վերացական մտածողություն, օրինաչափություններ գտնել։\par\medskip}
      \onslide<3->{Մարդը կարող է բաներ հասկանալ՝ \textbf{առանց ձեռք տալու}։\par}
    \end{column}
    \begin{column}{0.5\textwidth}
      \centering
      \begin{tikzpicture}[storybox/.style={flow=4.6cm, align=left, inner sep=4pt}]
        \node[storybox, visible on=<3->] (a) at (0,0) {\showon{3-}{\emoji{lock}}\ Դարպասը փակված է մեծ կողպեքներով};
        \node[storybox, visible on=<4->, fill=mint!40!cardwhite] (b) at (0,-1.4) {\showon{4-}{\emoji{thought}}\ «Դե, փակ է» $\rightarrow$ ուրիշ մուտք ես փնտրում};
        \node[storybox, visible on=<5->, fill=brainpink!30!cardwhite] (c) at (0,-2.8) {\showon{5-}{\emoji{muscle}}\ Առանց մտածելու քաշում ես դարպասը. «Հա, փակ է»};
        \draw[-{Stealth}, plum, line width=1pt, visible on=<4->] (a) -- (b);
      \end{tikzpicture}
    \end{column}
  \end{columns}
\end{frame}

% We can only measure it indirectly
\begin{frame}{Ինչպե՞ս չափել այն, ինչ չես տեսնում}
  \begin{columns}[c]
    \begin{column}{0.48\textwidth}
      \begin{itemize}
        \item Տեղեկության մշակումը ուղեղի ներսում ուղղակիորեն տեսնել չենք կարող
        \item Ամենաժամանակակից սկաներներն էլ, Բըրնեթի խոսքով, ցույց են տալիս «տարբեր գույների մշուշոտ բծեր»
        \item<2-> Մնում է չափել \textbf{անուղղակի}՝ վարքով ու հատուկ թեստերով
      \end{itemize}
    \end{column}
    \begin{column}{0.46\textwidth}
      \onslide<3->{%
      \begin{card}
        \emoji[1.3em]{judge}\ զանգվածը՝ կշռելով\par\smallskip
        \emoji[1.3em]{ruler}\ հասակը՝ չափելով\par\smallskip
        \emoji[1.3em]{police}\ ալկոհոլը՝ ոստիկանի սարքի մեջ փչելով\par\smallskip
        \emoji[1.3em]{numbers}\ խելքը՝ \textbf{IQ թեստով}
      \end{card}}
      \onslide<4->{\centering{\blackfont\large Պարզ է, չէ՞}\par}
      \onslide<5->{\centering\small ...ոչ այնքան\par}
    \end{column}
  \end{columns}
\end{frame}

% Myth or fact: the average IQ is always 100
\begin{frame}{Առասպել, թե՞ փաստ \#1}
  \mythorfact{Մարդկանց միջին IQ-ն միշտ 100 է}{ՓԱՍՏ}{mintdark}{%
    \begin{card}[colback=mint!35!cardwhite]
      IQ-ն (intelligence quotient՝ ինտելեկտի գործակից) կիլոգրամի նման չափում չէ, այլ ցույց է տալիս \textbf{քո տեղը մյուսների համեմատ}։ 100-ը միջինն է ըստ սահմանման։\par\smallskip
      Բըրնեթ. «X երկրում միջին IQ-ն ընդամենը 85 է» ասելը նույնն է, թե ասես՝ «X երկրում մետրը 85 սմ է»։
    \end{card}}
\end{frame}

% The bell curve with the book's bands
\begin{frame}{Զանգակաձև կոր}
  \centering
  % IQ is a normal distribution with mean 100 and SD 15 (Wechsler scales). Band names are the book's
% classification; the shares are computed from the normal curve (not in the book):
% <70 2.3%, 70-79 6.8%, 80-89 16.1%, 90-109 49.5%, 110-119 16.1%, 120-129 6.8%, 130+ 2.3%.
% No \foreach inside the axis: pgfplots draws nodes after the loop, so loop variables would be lost.
\begin{tikzpicture}
  \begin{axis}[
      width=13.2cm, height=4.0cm,
      xmin=45, xmax=155, ymin=0, ymax=0.031,
      axis x line*=bottom, axis y line=none, axis line style={plum},
      xtick={70,80,90,100,110,120,130}, tick label style={font=\small, text=plum},
      clip=false,
      declare function={g(\x)=exp(-((\x-100)^2)/450)/(15*sqrt(2*pi));},
    ]
    \addplot[draw=none, fill=amRed!70!cream, domain=45:70, samples=40] {g(x)} \closedcycle;
    \addplot[draw=none, fill=amRed!35!cream, domain=70:80, samples=20] {g(x)} \closedcycle;
    \addplot[draw=none, fill=amOrange!55!cream, domain=80:90, samples=20] {g(x)} \closedcycle;
    \addplot[draw=none, fill=amBlue!35!cream, domain=90:110, samples=40] {g(x)} \closedcycle;
    \addplot[draw=none, fill=amOrange!55!cream, domain=110:120, samples=20] {g(x)} \closedcycle;
    \addplot[draw=none, fill=amRed!35!cream, domain=120:130, samples=20] {g(x)} \closedcycle;
    \addplot[draw=none, fill=amRed!70!cream, domain=130:155, samples=40] {g(x)} \closedcycle;
    \addplot[plum, line width=1.2pt, domain=45:155, samples=160] {g(x)};
    \draw[plum!50, line width=0.5pt] (axis cs:70,0) -- (axis cs:70,0.0036);
    \draw[plum!50, line width=0.5pt] (axis cs:80,0) -- (axis cs:80,0.0109);
    \draw[plum!50, line width=0.5pt] (axis cs:90,0) -- (axis cs:90,0.0213);
    \draw[plum!50, line width=0.5pt] (axis cs:110,0) -- (axis cs:110,0.0213);
    \draw[plum!50, line width=0.5pt] (axis cs:120,0) -- (axis cs:120,0.0109);
    \draw[plum!50, line width=0.5pt] (axis cs:130,0) -- (axis cs:130,0.0036);
    % shares, on or above each band
    \node[font=\blackfont\large, text=plum] at (axis cs:100,0.012) {49.5\%};
    \node[font=\blackfont\footnotesize, text=plum] at (axis cs:85,0.0062) {16.1\%};
    \node[font=\blackfont\footnotesize, text=plum] at (axis cs:115,0.0062) {16.1\%};
    \node[font=\blackfont\footnotesize, text=plum] at (axis cs:75,0.0098) {6.8\%};
    \node[font=\blackfont\footnotesize, text=plum] at (axis cs:125,0.0098) {6.8\%};
    \node[font=\blackfont\footnotesize, text=plum] at (axis cs:61,0.0045) {2.3\%};
    \node[font=\blackfont\footnotesize, text=plum] at (axis cs:139,0.0045) {2.3\%};
    % the book's names for the bands, under the tick labels
    \node[anchor=north, font=\scriptsize, text=plum!85!cream, align=center] at (axis cs:58,-0.0042) {չափազանց\\ ցածր};
    \node[anchor=north, font=\scriptsize, text=plum!85!cream, align=center, yshift=-6.5mm] at (axis cs:75,-0.0042) {սահմանային};
    \node[anchor=north, font=\scriptsize, text=plum!85!cream, align=center] at (axis cs:85,-0.0042) {միջինից\\ ցածր};
    \node[anchor=north, font=\scriptsize\bfseries, text=plum, align=center] at (axis cs:100,-0.0042) {միջին};
    \node[anchor=north, font=\scriptsize, text=plum!85!cream, align=center] at (axis cs:115,-0.0042) {միջինից\\ բարձր};
    \node[anchor=north, font=\scriptsize, text=plum!85!cream, align=center, yshift=-6.5mm] at (axis cs:125,-0.0042) {բարձր};
    \node[anchor=north, font=\scriptsize, text=plum!85!cream, align=center] at (axis cs:142,-0.0042) {շատ\\ բարձր};
  \end{axis}
\end{tikzpicture}
\par\smallskip
  \begin{columns}[T]
    \begin{column}{0.48\textwidth}
      \begin{card}[colback=butter!45!cream]
        Թեստը չի չափում խելքի «հում» քանակը. այն ցույց է տալիս, թե \textbf{որտեղ ես կորի վրա}։
      \end{card}
    \end{column}
    \begin{column}{0.46\textwidth}
      \onslide<2->{%
      {\footnotesize\emoji{microscope}\ Գրքում՝ «80-ից 110-ի միջև է մարդկանց 80\%-ից ավելին»։ Իրականում այդտեղ է մոտ 66\%-ը, իսկ 80-ից 120-ի միջև՝ 82\%-ը։\par}}
    \end{column}
  \end{columns}
  \source{միջինը 100, ստանդարտ շեղումը 15 (Վեքսլերի սանդղակներ), տոկոսները հաշված են կորից}
\end{frame}

% Thought experiment: a virus that kills everyone above 100
\begin{frame}{Մտավոր փորձ. վիրուս, որ «ջնջում» է խելացիներին}
  \begin{columns}[T]
    \begin{column}{0.52\textwidth}
      \begin{enumerate}
        \item<1-> Վիրուսը վերացնում է բոլոր նրանց, ում IQ-ն 100-ից բարձր է
        \item<2-> Մնացածների միջին IQ-ն... նորից \textbf{100} է
        \item<3-> Երեկվա 99-ները դառնում են ինտելեկտուալ էլիտան \emoji{crown}
      \end{enumerate}
      \vspace{1mm}
      \onslide<4->{%
      \begin{card}[colback=butter!45!cream]
        \emoji{dollar}\ Դոլարը միշտ 100 ցենտ է, բայց թե ինչ կգնես դրանով, փոխվում է։ IQ-ն նույնպես. 100-ը միշտ միջինն է, բայց դրա «արժեքը» փոփոխական է։
      \end{card}}
    \end{column}
    \begin{column}{0.42\textwidth}
      \onslide<5->{%
      \begin{scienceupdate}
        Ավելի ճշգրիտ՝ 100-ը այն խմբի միջինն է, որի վրա թեստը \textbf{չափաբերվել} է (norming)։\par\smallskip
        Հին նորմերով այսօրվա մարդիկ 100-ից բարձր են ստանում, դրա համար թեստերը թարմացնում են (Ֆլինի էֆեկտ, գլխի վերջում)։
      \end{scienceupdate}}
    \end{column}
  \end{columns}
  \source{Flynn, Psychological Bulletin, 1987}
\end{frame}

% IQ 270 claims, Einstein, Guinness
\begin{frame}{IQ 270. Իսկապե՞ս}
  \begin{columns}[c]
    \begin{column}{0.58\textwidth}
      \begin{itemize}
        \item Այնշտայնի ու Հոքինգի IQ-ն, ըստ գնահատականների, մոտ \textbf{160} էր
        \item<2-> «Իմ IQ-ն 270 է» ասողը հավանաբար ոչ գիտական թեստ է անցել կամ սխալ է կարդացել արդյունքը
        \item<3-> Գինեսի գիրքը ժամանակին գրանցում էր 250-ից բարձր IQ-ներ, բայց \textbf{1990-ին} փակեց այդ անվանակարգը. այդ բարձրության վրա թեստերը շատ անորոշ են
      \end{itemize}
    \end{column}
    \begin{column}{0.36\textwidth}
      \centering
      \begin{tikzpicture}
        \node[draw=plum, line width=0.9pt, fill=cardwhite, rounded corners=3mm, inner sep=3mm, font=\blackfont\fontsize{30}{32}\selectfont, text=pinkdark] (n) {270};
        \node[anchor=north, font=\small, text=plum, align=center] at ([yshift=-2mm]n.south) {«ստուգել եմ\\ համացանցում»};
        \node[rotate=-12, visible on=<2->] at ([xshift=4mm, yshift=4mm]n.north east) {\showon{2-}{\bigemoji[1cm]{monocle}}};
      \end{tikzpicture}\par\medskip
      \onslide<3->{\bigemoji[1cm]{trophy}\ \raisebox{0.35cm}{\Large$\rightarrow$}\ \bigemoji[1cm]{wastebasket}}
    \end{column}
  \end{columns}
\end{frame}

% Real tests are instruments; extremes are hard to measure
\begin{frame}{Լոգարանի կշեռքը փղի համար}
  \begin{columns}[c]
    \begin{column}{0.55\textwidth}
      \begin{itemize}
        \item Գիտական IQ թեստերը իսկական \textbf{գործիքներ} են, ինչպես մանրադիտակը։ Թանկ են ու համացանցում \textbf{անվճար չեն բաժանվում}
        \item<2-> Ստեղծված են \textbf{միջին} խելքը չափելու համար. ծայրերում քիչ են օգտակար
        \item<3-> Ֆիզիկան դասարանում ցույց տալու համար բավական են կշռաքարերը, իսկ բարդ ֆիզիկայի համար պետք է մասնիկների արագացուցիչ
      \end{itemize}
    \end{column}
    \begin{column}{0.4\textwidth}
      \centering
      \onslide<4->{%
      \begin{tikzpicture}
        \filldraw[fill=cardwhite, draw=plum, line width=1pt, rounded corners=3mm] (-1.6,0) rectangle (1.6,0.7);
        \filldraw[fill=butter!60!cream, draw=plum, line width=0.8pt, rounded corners=1mm] (-0.55,0.15) rectangle (0.55,0.55);
        \node[font=\blackfont\scriptsize, text=pinkdark] at (0,0.35) {ERR};
        \draw[plum, line width=1.2pt] (-0.9,0.7) -- (-0.6,1.0) -- (-0.75,1.2) (0.5,0.7) -- (0.8,1.05);
        \draw[plum, line width=1pt, decorate, decoration={zigzag, amplitude=0.6mm, segment length=2mm}] (1.6,0.35) -- (2.3,0.75);
        \draw[plum, line width=1pt, decorate, decoration={zigzag, amplitude=0.6mm, segment length=2mm}] (-1.6,0.3) -- (-2.2,-0.1);
        \node[font=\small, text=plum] at (0,-0.35) {փիղը կշռելուց հետո};
      \end{tikzpicture}\par\medskip
      \begin{bookquote}
        Սովորական քաշերի համար լավն է, իսկ փղի դեպքում կստանաս միայն կոտրված պլաստմասսա ու զսպանակներ։
      \end{bookquote}}
    \end{column}
  \end{columns}
\end{frame}

% What a real test measures + verbal fluency try-it
\begin{frame}{Ինչ է ստուգում իսկական թեստը}
  \begin{columns}[T]
    \begin{column}{0.47\textwidth}
      \begin{card}
        \emoji[1.3em]{puzzle}\ \textbf{Տարածական ընկալում}՝ նախշերը լրացնելու խնդիրներով\par\smallskip
        \emoji[1.3em]{hourglass}\ \textbf{Ըմբռնման արագություն}՝ հատուկ հարցերով\par\smallskip
        \emoji[1.3em]{speech}\ \textbf{Խոսքի սահունություն} (verbal fluency)՝ որոշակի խմբի բառեր թվարկելով
      \end{card}
      \vspace{1mm}
      {\small Խելամիտ բաներ են, բայց գերհանճարի խելքի սահմանը դրանցով չես գտնի։\par}
    \end{column}
    \begin{column}{0.47\textwidth}
      \onslide<2->{%
      \begin{tryit}
        \textbf{30 վայրկյան}։ Գրի առ թղթի վրա որքան կարող ես շատ \textbf{կենդանիների անուններ}։\par\smallskip
        Հետո գրիր չաթում միայն թիվը։
      \end{tryit}}
      \onslide<3->{%
      \begin{card}[colback=lavender!35!cardwhite]
        Սա հենց խոսքի սահունության առաջադրանք է։ Մեկ փորձով ոչինչ չես իմանա քո մասին, դա պարզապես տաքացում է \emoji{grin}
      \end{card}}
    \end{column}
  \end{columns}
\end{frame}

% Pattern completion try-it
\begin{frame}{Փորձիր. ո՞րն է պակասող նկարը}
  \begin{columns}[c]
    \begin{column}{0.62\textwidth}
      \centering
      % A pattern-completion item in the style of the book's "spatial awareness with pattern completion tests":
% rows = shape, columns = count. Answer: Գ (three triangles), revealed on overlay 2.
\tikzset{
  pics/pzcirc/.style={code={\filldraw[fill=amBlue!55!cream, draw=plum, line width=0.6pt] (0,0) circle (0.12);}},
  pics/pzsq/.style={code={\filldraw[fill=amOrange!70!cream, draw=plum, line width=0.6pt] (-0.11,-0.11) rectangle (0.11,0.11);}},
  pics/pztri/.style={code={\filldraw[fill=amRed!55!cream, draw=plum, line width=0.6pt] (-0.13,-0.1) -- (0.13,-0.1) -- (0,0.14) -- cycle;}},
}
\begin{tikzpicture}
  % the 3x3 grid
  \foreach \r in {0,1,2}{\foreach \c in {0,1,2}{\node[ptile] at (\c*1.3,-\r*1.3) {};}}
  \foreach \r/\s in {0/pzcirc, 1/pzsq}{
    \foreach \c in {0,1,2}{\foreach \k in {0,...,\c}{\pic at ({\c*1.3+(\k-\c/2)*0.32},-\r*1.3) {\s};}}}
  \foreach \c in {0,1}{\foreach \k in {0,...,\c}{\pic at ({\c*1.3+(\k-\c/2)*0.32},-2.6) {pztri};}}
  \node[ptile, fill=butter!70!cream, font=\blackfont\LARGE, text=plum] at (2.6,-2.6) {?};
  % the four options
  \node[font=\blackfont\small, text=plum] at (5.55,0.75) {Ընտրիր մեկը};
  \foreach \i/\x/\y in {0/4.9/0, 1/6.2/0, 2/4.9/-1.75, 3/6.2/-1.75}{\node[ptile] (o\i) at (\x,\y) {};}
  \foreach \k in {0,1,2}{\pic at ({4.9+(\k-1)*0.32},0) {pzsq};}
  \foreach \k in {0,1}{\pic at ({6.2+(\k-0.5)*0.32},0) {pztri};}
  \foreach \k in {0,1,2}{\pic at ({4.9+(\k-1)*0.32},-1.75) {pztri};}
  \foreach \k in {0,1,2}{\pic at ({6.2+(\k-1)*0.32},-1.75) {pzcirc};}
  \foreach \i/\t in {0/Ա, 1/Բ, 2/Գ, 3/Դ}{\node[font=\blackfont\normalsize, text=plum, below=0.5mm of o\i] {\t};}
  % answer
  \draw[mintdark, line width=2pt, rounded corners=2mm, visible on=<2->] (4.25,-1.1) rectangle (5.55,-2.75);
\end{tikzpicture}

    \end{column}
    \begin{column}{0.34\textwidth}
      \begin{tryit}
        Գտիր օրինաչափությունը ու \textbf{գրիր տառը} չաթում։
      \end{tryit}
      \onslide<2->{%
      \begin{card}[colback=mint!35!cardwhite]
        \textbf{Գ}. տողում՝ նույն պատկերը, սյունակում՝ նույն քանակը։ Իսկական թեստերում խնդիրները շատ ավելի են բարդանում։
      \end{card}}
    \end{column}
  \end{columns}
\end{frame}

% The circularity problem and other criticisms
\begin{frame}{Շրջանաձև խնդիր}
  \begin{columns}[T]
    \begin{column}{0.5\textwidth}
      \begin{card}[colback=brainpink!25!cardwhite]
        \emoji[1.3em]{cyclone}\ Թեստերը չափում են խելքը, իսկ թե ինչ է խելքը, գիտենք... \textbf{թեստերից}։ Թերահավատ գիտնականներին սա դուր չի գալիս։
      \end{card}
      \onslide<3->{%
      \begin{card}[colback=lavender!35!cardwhite]
        Թեստերը բազմիցս վերանայվել ու ստուգվել են, բայց ոմանց կարծիքով բուն խնդիրը մնում է։
      \end{card}}
    \end{column}
    \begin{column}{0.44\textwidth}
      \onslide<2->{%
      {\blackfont Արդյունքի վրա ազդում են նաև}\par\smallskip
      \begin{itemize}
        \item դաստիարակությունն ու միջավայրը
        \item ընդհանուր առողջությունը
        \item թեստ հանձնելու հմտությունը
        \item կրթությունը
      \end{itemize}}
      \onslide<4->{\small Գուցե թեստերն օգտակար են, բայց \textbf{ոչ այն բանի համար}, ինչի համար ստեղծվել են։\par}
    \end{column}
  \end{columns}
\end{frame}

% Not all doom and gloom
\begin{frame}{Բայց ամեն ինչ այդքան մութ չէ}
  \begin{columns}[c]
    \begin{column}{0.54\textwidth}
      \begin{itemize}
        \item Այսօրվա թեստերը մեկ գնահատականի փոխարեն տալիս են \textbf{շատ ենթաթեստեր}՝ տարածական ընկալում, թվաբանություն և այլն
        \item<2-> Մարդու արդյունքը կյանքի ընթացքում \textbf{բավական կայուն} է մնում
        \item<3-> Ուրեմն թեստերը բռնում են ինչ-որ \textbf{ներքին} հատկություն, ոչ թե պատահական հանգամանքներ
      \end{itemize}
    \end{column}
    \begin{column}{0.4\textwidth}
      \centering
      \onslide<4->{%
      \mascot[1.6cm]{happy}\par\smallskip
      \begin{bookquote}
        Խելքի նշաններից մեկը սեփական չիմացածի գիտակցումն է։ Ապրե՛ս։
      \end{bookquote}}
      \vspace{5mm}
    \end{column}
  \end{columns}
  \source{Nisbett et al., American Psychologist, 2012}
\end{frame}

% Meme: Expanding Brain, IQ claims
\begin{frame}{Որքա՞ն է քո IQ-ն}
  \centering
  \begin{meme}[13.4cm]{expanding_brain_row.jpg}
    \memecap[text width=2.9cm, font=\blackfont\small]{0.121}{0.18}{Իմ IQ-ն 120 է}
    \memecap[text width=2.9cm, font=\blackfont\small]{0.367}{0.18}{Իմ IQ-ն 160 է,\\ Այնշտայնի պես}
    \memecap[text width=2.9cm, font=\blackfont\small]{0.627}{0.18}{Իմ IQ-ն 270 է.\\ ինստագրամն ասաց}
    \memecap[text width=2.9cm, font=\blackfont\small]{0.883}{0.18}{Չգիտեմ՝ որքան է,\\ ու դա նորմալ է}
  \end{meme}
\end{frame}
